\documentclass{harryopo-paper}

\title{基于注意力机制的轻量化图像分类方法研究}
\author{张三 计算机学院 2025000101\\ \fzfs\small 示例大学}
\date{2026年08月30日}


\begin{document}

\maketitle
\begin{abstract}
图像分类是计算机视觉领域的基础任务之一。本文提出了一种基于注意力机制的轻量化图像分类方法，在保持较高准确率的同时显著降低了模型复杂度。实验结果表明，所提方法在 ImageNet 数据集上的 Top-1 准确率达到 78.3\%，参数量仅为 3.2M，FLOPs 为 0.8G，优于现有的 MobileNetV3 和 ShuffleNetV2 等轻量化模型。
\end{abstract}
\keywords{图像分类；注意力机制；轻量化网络；深度学习}

\subsection*{一、引言}
\addcontentsline{toc}{subsection}{一、引言}

图像分类是计算机视觉中最基本的任务之一，其目标是将输入图像分配到预定义的类别中。传统方法依赖手工设计的特征提取器，泛化能力有限。深度学习技术的兴起——特别是卷积神经网络（CNN）——彻底改变了这一局面。

自 AlexNet 在 2012 年 ImageNet 挑战赛中取得突破以来，研究者们提出了大量 CNN 架构。ResNet 通过残差连接解决了深层网络训练困难的问题，EfficientNet 则通过复合缩放策略实现了效率与精度的平衡。

\subsubsection*{1.1 研究背景}
\addcontentsline{toc}{subsubsection}{1.1 研究背景}

随着移动设备和边缘计算的普及，模型轻量化成为研究热点。轻量化模型需要在保持较高准确率的同时，显著减少参数量和计算量。

\subsubsection*{1.2 主要贡献}
\addcontentsline{toc}{subsubsection}{1.2 主要贡献}

本文的主要贡献包括：

\begin{enumerate}
  \item 提出了一种新的通道注意力机制，{\fzht 计算开销降低 40\%}
  \item 设计了自适应空间池化模块，增强多尺度特征提取能力
  \item 在 ImageNet、CIFAR-100 上均取得 SOTA 性能
\end{enumerate}

\subsection*{二、相关工作}
\addcontentsline{toc}{subsection}{二、相关工作}

\subsubsection*{2.1 轻量化网络}
\addcontentsline{toc}{subsubsection}{2.1 轻量化网络}

MobileNetV3 通过深度可分离卷积和 SE 注意力机制实现高效推理；ShuffleNetV2 则通过通道混洗优化并行计算。

\begin{table}[htbp]
  \centering
  \small
  \begin{tabularx}{\textwidth}{>{\raggedright\arraybackslash}X >{\raggedright\arraybackslash}X >{\raggedright\arraybackslash}X >{\raggedright\arraybackslash}X}
    \toprule
    {\fzht 模型} & {\fzht Top-1 (\%)} & {\fzht Params (M)} & {\fzht FLOPs (G)} \\
    \midrule
    MobileNetV3-Small & 67.5 & 2.5 & 0.06 \\
    ShuffleNetV2-1.0× & 69.4 & 2.3 & 0.15 \\
    EfficientNet-B0 & 77.3 & 5.3 & 0.39 \\
    {\fzht 本文方法} & {\fzht 78.3} & {\fzht 3.2} & {\fzht 0.8} \\
    \bottomrule
  \end{tabularx}
\end{table}

\subsubsection*{2.2 注意力机制}
\addcontentsline{toc}{subsubsection}{2.2 注意力机制}

SE 注意力通过全局平均池化建模通道关系，CBAM 进一步引入空间注意力。两者均能以较小开销提升模型性能。

\subsection*{三、方法}
\addcontentsline{toc}{subsection}{三、方法}

\subsubsection*{3.1 整体架构}
\addcontentsline{toc}{subsubsection}{3.1 整体架构}

本文方法的整体架构如下图所示：

\begin{lstlisting}[style=plainstyle,caption={}]
输入图像 → 初始卷积 → 阶段1（注意力模块 × 3）→ 阶段2（注意力模块 × 4）
       → 阶段3（注意力模块 × 6）→ 阶段4（注意力模块 × 3）→ 全局池化 → 分类器
\end{lstlisting}

\subsubsection*{3.2 轻量化注意力模块}
\addcontentsline{toc}{subsubsection}{3.2 轻量化注意力模块}

轻量化注意力模块（LAM）由两个分支组成：通道注意力分支和空间注意力分支。

\begin{quote}
{\fzht 公式 1} 通道注意力计算公式：
\$$M_c(F) = \sigma(\text{MLP}(\text{AvgPool}(F)) + \text{MLP}(\text{MaxPool}(F)))$\$
\end{quote}

其中 $\sigma$ 表示 Sigmoid 激活函数，$\text{MLP}$ 为多层感知机。

\begin{lstlisting}[style=pystyle,caption={}]
class LightweightAttention(nn.Module):
    def __init__(self, channels, reduction=16):
        super().__init__()
        self.avg_pool = nn.AdaptiveAvgPool2d(1)
        self.max_pool = nn.AdaptiveMaxPool2d(1)
        self.mlp = nn.Sequential(
            nn.Linear(channels, channels // reduction),
            nn.ReLU(inplace=True),
            nn.Linear(channels // reduction, channels),
            nn.Sigmoid()
        )

    def forward(self, x):
        b, c, _, _ = x.size()
        avg_out = self.mlp(self.avg_pool(x).view(b, c))
        max_out = self.mlp(self.max_pool(x).view(b, c))
        scale = (avg_out + max_out).view(b, c, 1, 1)
        return x * scale.expand_as(x)
\end{lstlisting}

\subsubsection*{3.3 损失函数}
\addcontentsline{toc}{subsubsection}{3.3 损失函数}

本文使用标签平滑交叉熵损失：

\begin{quote}
{\fzht 公式 2} 标签平滑交叉熵损失：
\$$L = -\sum_{i=1}^{C} y_i^{\text{smooth}} \log(p_i), \quad y_i^{\text{smooth}} = (1-\epsilon) y_i + \frac{\epsilon}{C}$\$
\end{quote}

其中 $\epsilon = 0.1$ 为平滑系数，$C$ 为类别数。

\subsection*{四、实验与分析}
\addcontentsline{toc}{subsection}{四、实验与分析}

\subsubsection*{4.1 实验设置}
\addcontentsline{toc}{subsubsection}{4.1 实验设置}

\begin{itemize}
  \item 数据集：ImageNet-1K（1000 类，128 万张训练图像）
  \item 优化器：SGD（momentum=0.9, weight\_decay=4e-5）
  \item 初始学习率：0.1，cosine 衰减，训练 300 epoch
  \item 批量大小：256（4×RTX 3090）
\end{itemize}

\subsubsection*{4.2 消融实验}
\addcontentsline{toc}{subsubsection}{4.2 消融实验}

\begin{table}[htbp]
  \centering
  \small
  \begin{tabularx}{\textwidth}{>{\raggedright\arraybackslash}X >{\raggedright\arraybackslash}X >{\raggedright\arraybackslash}X >{\raggedright\arraybackslash}X}
    \toprule
    {\fzht 模块} & {\fzht Top-1 (\%)} & {\fzht Params (M)} & {\fzht FLOPs (G)} \\
    \midrule
    基线 & 72.4 & 2.8 & 0.5 \\
    + 通道注意力 & 76.1 & 3.0 & 0.7 \\
    + 空间注意力 & 77.5 & 3.1 & 0.75 \\
    + 标签平滑 & 78.3 & 3.2 & 0.8 \\
    \bottomrule
  \end{tabularx}
\end{table}

\begin{quote}
{\fzht 结论}：三个模块均带来稳定提升，组合使用时效果最佳。
\end{quote}

\subsubsection*{4.3 可视化分析}
\addcontentsline{toc}{subsubsection}{4.3 可视化分析}

通过 Grad-CAM 可视化发现，注意力模块能有效聚焦于图像的判别性区域，抑制背景噪声的干扰。

\subsection*{五、结论}
\addcontentsline{toc}{subsection}{五、结论}

本文提出了一种基于注意力机制的轻量化图像分类方法，通过通道-空间双分支注意力设计，在仅增加 14\% 参数量的前提下将 Top-1 准确率从 72.4\% 提升至 78.3\%。实验结果表明所提方法在精度和效率间取得了良好平衡，适合在边缘设备上部署。

未来的研究方向包括：

\begin{enumerate}
  \item 探索神经架构搜索（NAS）自动设计注意力模块
  \item 研究量化感知训练，进一步压缩模型体积
  \item 拓展到目标检测、语义分割等下游任务
\end{enumerate}

\begin{thebibliography}{99}
\bibitem{ref1} Krizhevsky A, Sutskever I, Hinton G E. ImageNet classification with deep convolutional neural networks[C]. NeurIPS, 2012: 1097-1105.
\bibitem{ref2} He K, Zhang X, Ren S, et al. Deep residual learning for image recognition[C]. CVPR, 2016: 770-778.
\bibitem{ref3} Tan M, Le Q. EfficientNet: Rethinking model scaling for convolutional neural networks[C]. ICML, 2019: 6105-6114.
\bibitem{ref4} Howard A, Sandler M, Chu G, et al. Searching for MobileNetV3[C]. ICCV, 2019: 1314-1324.
\bibitem{ref5} Ma N, Zhang X, Zheng H T, et al. ShuffleNetV2: Practical guidelines for efficient CNN architecture design[C]. ECCV, 2018: 116-131.
\end{thebibliography}

\end{document}